    Solutions for the second day problems at the IMC 2000

    Problem 1.
    a) Show that the unit square can be partitioned into n smaller squares if n is large
enough.
    b) Let d ≥ 2. Show that there is a constant N (d) such that, whenever n ≥ N (d), a
d-dimensional unit cube can be partitioned into n smaller cubes.

     Solution. We start with the following lemma: If a and b be coprime positive integers
then every sufficiently large positive integer m can be expressed in the form ax + by with
x, y non-negative integers.
     Proof of the lemma. The numbers 0, a, 2a, . . . , (b−1)a give a complete residue system
modulo b. Consequently, for any m there exists a 0 ≤ x ≤ b − 1 so that ax ≡ m (mod b).
If m ≥ (b − 1)a, then y = (m − ax)/b, for which x + by = m, is a non-negative integer, too.
     Now observe that any dissection of a cube into n smaller cubes may be refined to
give a dissection into n + (ad − 1) cubes, for any a ≥ 1. This refinement is achieved by
picking an arbitrary cube in the dissection, and cutting it into ad smaller cubes. To prove
the required result, then, it suffices to exhibit two relatively prime integers of form a d − 1.
In the 2-dimensional case, a1 = 2 and a2 = 3 give the coprime numbers 22 − 1 = 3 and
32 − 1 = 8. In the general case, two such integers are 2d − 1 and (2d − 1)d − 1, as is easy
to check.

     Problem 2. Let f be continuous and nowhere monotone on [0, 1]. Show that the set
of points on which f attains local minima is dense in [0, 1].
     (A function is nowhere monotone if there exists no interval where the function is
monotone. A set is dense if each non-empty open interval contains at least one element of
the set.)

     Solution. Let (x − α, x + α) ⊂ [0, 1] be an arbitrary non-empty open interval. The
function f is not monoton in the intervals [x − α, x] and [x, x + α], thus there exist some
real numbers x − α ≤ p < q ≤ x, x ≤ r < s ≤ x + α so that f (p) > f (q) and f (r) < f (s).
     By Weierstrass’ theorem, f has a global minimum in the interval [p, s]. The values f (p)
and f (s) are not the minimum, because they are greater than f (q) and f (s), respectively.
Thus the minimum is in the interior of the interval, it is a local minimum. So each non-
empty interval (x − α, x + α) ⊂ [0, 1] contains at least one local minimum.

     Problem 3. Let p(z) be a polynomial of degree n with complex coefficients. Prove
that there exist at least n + 1 complex numbers z for which p(z) is 0 or 1.

     Solution. The statement is not true if p is a constant polynomial. We prove it only
in the case if n is positive.
     For an arbitrary polynomial q(z) and complex number c, denote by µ(q, c) the largest
exponent α for which q(z) is divisible by (z − c)α . (With other words, if c is a root of q,
then µ(q, c) is the root’s multiplicity. Otherwise 0.)

                                               1
     Denote by S0 and S1 the sets of complex numbers z for which p(z) is 0 or 1, respec-
tively. These sets contain all roots of the polynomials p(z) and p(z) − 1, thus
                                      X                 X
                                            µ(p, c) =          µ(p − 1, c) = n.                          (1)
                                     c∈S0               c∈S1


The polynomial p0 has at most n − 1 roots (n > 0 is used here). This implies that
                                              X
                                                       µ(p0 , c) ≤ n − 1.                                (2)
                                            c∈S0 ∪S1


    If p(c) = 0 or p(c) − 1 = 0, then

                        µ(p, c) − µ(p0 c) = 1 or µ(p − 1, c) − µ(p0 c) = 1,                              (3)

respectively. Putting (1), (2) and (3) together we obtain
                               X                           X                        
             S0 + S1 =                 µ(p, c) − µ(p0 , c) +  µ(p − 1, c) − µ(p0 , c) =
                               c∈S0                                c∈S1

           X                  X                          X
       =          µ(p, c) +          µ(p − 1, c) −                µ(p0 , c) ≥ n + n − (n − 1) = n + 1.
           c∈S0               c∈S1                     c∈S0 ∪S1



     Problem 4. Suppose the graph of a polynomial of degree 6 is tangent to a straight
line at 3 points A1 , A2 , A3 , where A2 lies between A1 and A3 .
     a) Prove that if the lengths of the segments A1 A2 and A2 A3 are equal, then the areas
of the figures bounded by these segments and the graph of the polynomial are equal as well.
                 A2 A3
     b) Let k =        , and let K be the ratio of the areas of the appropriate figures. Prove
                 A1 A2
that
                                       2 5          7
                                         k < K < k5 .
                                       7            2

     Solution. a) Without loss of generality, we can assume that the point A2 is the origin
of system of coordinates. Then the polynomial can be presented in the form

                          y = a0 x4 + a1 x3 + a2 x2 + a3 x + a4 x2 + a5 x,
                                                               


where the equation y = a5 x determines the straight line A1 A3 . The abscissas of the points
A1 and A3 are −a and a, a > 0, respectively. Since −a and a are points of tangency, the
numbers −a and a must be double roots of the polynomial a0 x4 + a1 x3 + a2 x2 + a3 x + a4 .
It follows that the polynomial is of the form

                                            y = a0 (x2 − a2 )2 + a5 x.

                                                           2
      The equality follows from the equality of the integrals

                            Z0                         Za
                                 a0 x − a x2 dx =
                                    2    2
                                                            a0 x2 − a2 x2 dx
                                                                     

                         −a                            0


due to the fact that the function y = a0 (x2 − a2 ) is even.
    b) Without loss of generality, we can assume that a0 = 1. Then the function is of the
form
                               y = (x + a)2 (x − b)2 x2 + a5 x,
where a and b are positive numbers and b = ka, 0 < k < ∞. The areas of the figures at
the segments A1 A2 and A2 A3 are equal respectively to

                        Z0
                                                             a7
                             (x + a)2 (x − b)2 x2 dx =          (7k 2 + 7k + 2)
                                                            210
                       −a


and
                        Zb
                                                             a7
                             (x + a)2 (x − b)2 x2 dx =          (2k 2 + 7k + 7)
                                                            210
                        0

Then
                                                 2k 2 + 7k + 7
                                        K = k5                 .
                                                 7k 2 + 7k + 2
                                             2
The derivative of the function f (k) = 2k  +7k+7
                                       7k2 +7k+2 is negative for 0 < k < ∞. Therefore f (k)
                                                                                  2
decreases from 27 to 27 when k increases from 0 to ∞. Inequalities 27 < 7k
                                                                        2k +7k+7     7
                                                                           2 +7k+2 < 2 imply

the desired inequalities.

    Problem 5. Let R+ be the set of positive real numbers. Find all functions f : R+ →
R+ such that for all x, y ∈ R+

                                     f (x)f (yf (x)) = f (x + y).

     First solution. First, if we assume that f (x) > 1 for some x ∈ R+ , setting y =
    x
           gives the contradiction f (x) = 1. Hence f (x) ≤ 1 for each x ∈ R+ , which implies
f (x) − 1
that f is a decreasing function.
     If f (x) = 1 for some x ∈ R+ , then f (x + y) = f (y) for each y ∈ R+ , and by the
monotonicity of f it follows that f ≡ 1.
     Let now f (x) < 1 for each x ∈ R+ . Then f is strictly decreasing function, in particular
injective. By the equalities

                                    f (x)f (yf (x)) = f (x + y) =

                                                   3
                                                                             
        = f yf (x) + x + y(1 − f (x)) = f (yf (x))f x + y(1 − f (x)) f (yf (x))

                                                                                             1 − f (1)
we obtain that x = (x + y(1 − f (x)))f (yf (x)). Setting x = 1, z = xf (1) and a =                     ,
                                                                                               f (1)
                    1
we get f (z) =          .
                 1 + az
                                                              1
    Combining the two cases, we conclude that f (x) =              for each x ∈ R+ , where
                                                           1 + ax
a ≥ 0. Conversely, a direct verification shows that the functions of this form satisfy the
initial equality.
     Second solution. As in the first solution we get that f is a decreasing function, in
particular differentiable almost everywhere. Write the initial equality in the form

                            f (x + y) − f (x)           f (yf (x)) − 1
                                              = f 2 (x)                .
                                    y                       yf (x)

It follows that if f is differentiable at the point x ∈ R+ , then there exists the limit
       f (z) − 1                     0           2                     +
                                                                                1 0
 lim             =: −a. Therefore f (x) = −af (x) for each x ∈ R , i.e.                 = a,
z→0+       z                                                                    f (x)
                               1
which means that f (x) =           . Substituting in the initial relaton, we find that b = 1
                            ax + b
and a ≥ 0.
                                                                           ∞
                                                                             1
    Problem 6. For an m × m real matrix A, eA is defined as                           An . (The sum is
                                                                           P
                                                                                 n!
                                                                           n=0
convergent for all matrices.) Prove or disprove, that for all real polynomials p and m × m
real matrices A and B, p(eAB ) is nilpotent if and only if p(eBA ) is nilpotent. (A matrix
A is nilpotent if Ak = 0 for some positive integer k.)
     Solution. First we prove that for any polynomial q and m × m matrices A and B,
the characteristic polinomials of q(eAB ) and q(eBA ) are the same. It is easy to check that
                             ∞
for any matrix X, q(eX ) =      cn X n with some real numbers cn which depend on q. Let
                             P
                              n=0

                             ∞
                             X                        ∞
                                                      X
                        C=         cn · (BA)n−1 B =         cn · B(AB)n−1 .
                             n=1                      n=1


Then q(eAB ) = c0 I + AC and q(eBA ) = c0 I + CA. It is well-known that the characteristic
polynomials of AC and CA are the same; denote this polynomial by f (x). Then the
characteristic polynomials of matrices q(eAB ) and q(eBA ) are both f (x − c0 ).
                                                                  k
     Now assume that the matrix p(eAB ) is nilpotent, i.e. p(eAB ) = 0 for some positive
integer k. Chose q = pk . The characteristic polynomial of the matrix q(eAB ) = 0 is xm ,
so the same holds for the matrix q(eBA ). By the theorem of Cayley and Hamilton, this
                     m           km
implies that q(eBA ) = p(eBA )         = 0. Thus the matrix q(eBA ) is nilpotent, too.



                                                 4
